\chapter{Event Selection}
\label{sec:eventSel}

Candidate \Wln{} events are collected in data with single-lepton triggers. The primary, unprescaled triggers used in both 5 and 13~\TeV\ datasets are the $\texttt{HLT\_e15\_lhloose\_nod0\_L1EM12}$ and $\texttt{HLT\_mu14}$. These single lepton triggers require either at least one reconstructed electron with \ET larger than 15~\GeV\ passing \texttt{loose} identification requirements or one muon with \pt larger than 14~\GeV.

Selected events are required to contain exactly one identified and isolated lepton (muon or electron) candidate satisfying the criteria described in Section~\ref{sec:objects}.
Events with additional leptons of the same flavour with transverse momentum greater than 20 ~\GeV{} are discarded to lower the \Zboson\ background: the identification for those veto leptons point is \texttt{medium} for the muon channel and \texttt{loose} for the electron channel. There is no requirement on the number of leptons with different flavour than the channel under study.

To reduce background contamination, in particular from multijet
processes, events are required to have \MET{} , defined in
Eq.~(\ref{eq:METandNeutrino}), greater than 25~\GeV{}.  The \Wboson{}
boson transverse mass \mt{} is required to be larger than 50~\GeV,
where \mt{} is defined as usual:
\begin{equation}\label{eq:mt}
{m}_{T} = \sqrt{2 {p}_{T}^{\nu}p_{T} ^{l} (1-\cos\Delta\phi^{\nu})}
\end{equation}
and as in Eq.~(\ref{eq:METandNeutrino}), the neutrino $\pt^\nu$ and $\phi^\nu$ are estimated by the corresponding \MET{} quantities. The \ut criteria is defined as $\ut < 25 \GeV$.

In
Tables \todo{Need all the tables with the event selection cutflows}
$W^{\pm} \rightarrow \ell^{\pm}\nu$ signal selection event yields for
the $\sqrt{s} = 5~\TeV$ dataset are shown, and the corresponding
numbers for the 13~\TeV\ dataset are shown in
Tables \todo{add tables}. A
summary of observed compared to expected yields is then provided in
Table \todo{add tables}.  The yields for the signal and for the
electroweak and top backgrounds are obtained from MC samples.  The
events referred to as ``\Wln{} BG'' in the tables and plots refer to
the sum of \Wtaunu\ events and \Wln\ events of the wrong generated
charge passing due to charge-misidentification.  The estimate of the
multijet background passing due selection of leptons produced in
semi-leptonic decays of heavy quarks, in-flight decays, or
misidentification of photon conversions or hadrons, is done performed with
data-driven techniques described in Section~\ref{sec:bkgs}.


